\section{Solutions for Chapter~\ref{ch:vertex_algebras}}

\paragraph{Exercise~\ref{ex:voa-weight}.}
\phantomsection
\label{sol:ex:voa-weight}
The singular products have $n=0,\ldots,6$ and respective weights
$6,5,4,3,2,1,0$. Products with $n>6$ vanish because their
weights would be negative. CFT type gives $V_0=\C\one$,
so $a_{(6)}b$ is a scalar multiple of the vacuum; $a_{(4)}b$
is the weight-two candidate. This list permits a product to be
zero; weight counting does not establish its nonvanishing.



\paragraph{Exercise~\ref{ex:voa-cocycle}.}
\phantomsection
\label{sol:ex:voa-cocycle}
Write the exponent of the sign as a bilinear form $B(a,b)$
over $\Z/2\Z$. Bilinearity gives
\[
 B(a,b)+B(a+b,c)=B(b,c)+B(a,b+c),
\]
which is exactly the two-cocycle identity. For the triangular
root-basis choice, $B(a,b)-B(b,a)$ is
$\sum_{i>j}G_{ij}(a_ib_j+a_jb_i)$ modulo two.
The diagonal terms in $(a,b)$ are even, so this is $(a,b)$
modulo two. Hence the commutator ratio is $(-1)^{(a,b)}$.
A coboundary has the form
$f(a)f(b)/f(a+b)$ and is symmetric in $a,b$; it cancels
out of that ratio.



\paragraph{Exercise~\ref{ex:voa-rankone}.}
\phantomsection
\label{sol:ex:voa-rankone}
The oscillator series begins $1+q+2q^2+\cdots$.
For the vacuum lattice, $\Theta_L=1+2q+O(q^4)$, giving
\[
 \Theta_L\prod_{n\ge1}(1-q^n)^{-1}=1+3q+4q^2+\cdots.
\]
For the shifted lattice, the momentum weights are
$(m+1/2)^2$. Thus
\[
 \Theta_{L+\alpha/2}=q^{1/4}(2+2q^2+\cdots),\qquad
 \frac{\Theta_{L+\alpha/2}}{\prod_{n\ge1}(1-q^n)}
 =q^{1/4}(2+2q+6q^2+\cdots).
\]
The two shortest momenta explain the two-dimensional top.
The module is nevertheless invertible: adding its discriminant
coset twice gives $L$. An invertible simple has categorical
dimension one, independently of its top multiplicity.



\paragraph{Exercise~\ref{ex:voa-heisenberg-c2}.}
\phantomsection
\label{sol:ex:voa-heisenberg-c2}
Every oscillator of index $n\ge2$ can be removed modulo $C_2$:
\[
 J_{-n}b=\frac1{(n-1)!}(\mathsf T^{n-2}J)_{(-2)}b
 \in C_2(V).
\]
Commuting the creation oscillators shows that only powers of
$J_{-1}\one$ remain. To see that these powers have no relation,
count oscillators. A monomial of weight $w$ has at most $w$
oscillators, with equality only when every index is one.
For homogeneous Fock monomials $a,b$, the state $a_{(-2)}b$
has weight $\wt(a)+\wt(b)+1$, but Wick products have at most
$\wt(a)+\wt(b)$ oscillators; contractions only lower their number.
Thus every resulting monomial contains an index at least two.
No $C_2$ term has a pure $J_{-1}$ component in the Fock basis.
The quotient is $\C[x]$, an infinite-dimensional algebra.
The VOA is strongly generated by $J$, so this is the desired
counterexample to an implication from finite strong generation.



\paragraph{Exercise~\ref{ex:voa-coset}.}
\phantomsection
\label{sol:ex:voa-coset}
Subtract the charge-one Cartan boson:
\[
 \frac{3k}{k+2}-1=\frac{2(k-1)}{k+2}.
\]
At $k=3$ and $13$ this is $4/5$ and $8/5$.
Identification with the parafermion factors additionally needs
their generating fields, OPEs, null relations, and the proof
that the relevant simple quotient is realized. Central charge
is only one invariant of a VOA.



\paragraph{Exercise~\ref{ex:voa-gauge}.}
\phantomsection
\label{sol:ex:voa-gauge}
The transported basis squared norms are $4,40,40$, since $g^{-1}$
rescales each vector by the inverse of the stated factor.
Each of the last two coefficients has squared modulus
$(13+27)/400=1/10$. Therefore
\[
 \norm v_g^2=4+40/10+40/10=12.
\]
Adding coefficient magnitudes alone would give $6/5$,
which uses the old orthonormal form on the new coordinates.

